\documentclass[10pt,a4paper]{amsart}
\usepackage{amsmath,amssymb,amsthm,graphicx}
% Retained preamble for the Stacks Project chapter divisors.tex (Divisors).
% The chapter file pinned-original-source.tex carries only \input{preamble} on its first line, so no
% preamble text lives inside the chapter itself.  The Stacks Project uses one single shared preamble for
% every chapter; the snapshot stages it verbatim in sources/stacks-chow/preamble.tex.  The block below is
% copied mechanically, line for line, from that file's lines 14-18 (the Xy-pic setup) and lines 154-217
% (the native theorem declarations together with the mathematical macros).  Nothing was added, removed or
% redefined.  The wrapper class and the amsmath/amssymb/amsthm/graphicx base packages are supplied by the
% assembler, so the class/\externaldocument part of the shared preamble is not repeated here.
% For commutative diagrams we use Xy-pic
\usepackage[all]{xy}

% We use 2cell for 2-commutative diagrams.
\xyoption{2cell}
% Theorem environments.
%
\theoremstyle{plain}
\newtheorem{theorem}[subsection]{Theorem}
\newtheorem{proposition}[subsection]{Proposition}
\newtheorem{lemma}[subsection]{Lemma}

\theoremstyle{definition}
\newtheorem{definition}[subsection]{Definition}
\newtheorem{example}[subsection]{Example}
\newtheorem{exercise}[subsection]{Exercise}
\newtheorem{situation}[subsection]{Situation}

\theoremstyle{remark}
\newtheorem{remark}[subsection]{Remark}
\newtheorem{remarks}[subsection]{Remarks}

\numberwithin{equation}{subsection}

% Macros
%
\def\lim{\mathop{\mathrm{lim}}\nolimits}
\def\colim{\mathop{\mathrm{colim}}\nolimits}
\def\Spec{\mathop{\mathrm{Spec}}}
\def\Hom{\mathop{\mathrm{Hom}}\nolimits}
\def\Ext{\mathop{\mathrm{Ext}}\nolimits}
\def\SheafHom{\mathop{\mathcal{H}\!\mathit{om}}\nolimits}
\def\SheafExt{\mathop{\mathcal{E}\!\mathit{xt}}\nolimits}
\def\Sch{\mathit{Sch}}
\def\Mor{\mathop{\mathrm{Mor}}\nolimits}
\def\Ob{\mathop{\mathrm{Ob}}\nolimits}
\def\Sh{\mathop{\mathit{Sh}}\nolimits}
\def\NL{\mathop{N\!L}\nolimits}
\def\CH{\mathop{\mathrm{CH}}\nolimits}
\def\proetale{{pro\text{-}\acute{e}tale}}
\def\etale{{\acute{e}tale}}
\def\QCoh{\mathit{QCoh}}
\def\Ker{\mathop{\mathrm{Ker}}}
\def\Im{\mathop{\mathrm{Im}}}
\def\Coker{\mathop{\mathrm{Coker}}}
\def\Coim{\mathop{\mathrm{Coim}}}

% Boxtimes
%
\DeclareMathSymbol{\boxtimes}{\mathbin}{AMSa}{"02}

%
% Macros for moduli stacks/spaces
%
\def\QCohstack{\mathcal{QC}\!\mathit{oh}}
\def\Cohstack{\mathcal{C}\!\mathit{oh}}
\def\Spacesstack{\mathcal{S}\!\mathit{paces}}
\def\Quotfunctor{\mathrm{Quot}}
\def\Hilbfunctor{\mathrm{Hilb}}
\def\Curvesstack{\mathcal{C}\!\mathit{urves}}
\def\Polarizedstack{\mathcal{P}\!\mathit{olarized}}
\def\Complexesstack{\mathcal{C}\!\mathit{omplexes}}
% \Pic is the operator that assigns to X its picard group, usage \Pic(X)
% \Picardstack_{X/B} denotes the Picard stack of X over B
% \Picardfunctor_{X/B} denotes the Picard functor of X over B
\def\Pic{\mathop{\mathrm{Pic}}\nolimits}
\def\Picardstack{\mathcal{P}\!\mathit{ic}}
\def\Picardfunctor{\mathrm{Pic}}
\def\Deformationcategory{\mathcal{D}\!\mathit{ef}}

% --- appended 2026-10-12 (Chief): the project own definitions of the reference/slogan/history
% display environments, copied mechanically from the shared Stacks preamble lines 8-12
% (\usepackage{verbatim} plus the three \newenvironment lines, which the chapter preamble keeps
% verbatim from the project preamble). Nothing added, removed or redefined: these environments
% carry no mathematical content and the project itself typesets them as comments. Without these
% three lines a counted statement whose own text uses them fails with Environment slogan undefined.
\usepackage{verbatim}
\newenvironment{reference}{\comment}{\endcomment}
\newenvironment{slogan}{\comment}{\endcomment}
\newenvironment{history}{\comment}{\endcomment}
\usepackage{cleveref}
\crefname{originallemma}{original lemma}{original lemmas}
\Crefname{originallemma}{Original lemma}{Original lemmas}
\makeatletter
\newlabel{algebra-lemma-S-plus-generated}{{07Z4}{source}{}{sourcepointer-1}{}}
\newlabel{algebra-lemma-S-plus-generated@cref}{{[originallemma][0][]07Z4}{[1][0][]source}{}{}{}}
\makeatother
\crefname{originallemma}{original lemma}{original lemmas}
\Crefname{originallemma}{Original lemma}{Original lemmas}
\makeatletter
\newlabel{constructions-lemma-proj-quasi-compact}{{01MD}{source}{}{sourcepointer-2}{}}
\newlabel{constructions-lemma-proj-quasi-compact@cref}{{[originallemma][0][]01MD}{[1][0][]source}{}{}{}}
\makeatother
\crefname{originalsection}{original section}{original sections}
\Crefname{originalsection}{Original section}{Original sections}
\makeatletter
\newlabel{constructions-section-relative-proj-via-glueing}{{01NM}{source}{}{sourcepointer-3}{}}
\newlabel{constructions-section-relative-proj-via-glueing@cref}{{[originalsection][0][]01NM}{[1][0][]source}{}{}{}}
\makeatother
\crefname{originallemma}{original lemma}{original lemmas}
\Crefname{originallemma}{Original lemma}{Original lemmas}
\makeatletter
\newlabel{schemes-lemma-quasi-compact-affine}{{01K4}{source}{}{sourcepointer-4}{}}
\newlabel{schemes-lemma-quasi-compact-affine@cref}{{[originallemma][0][]01K4}{[1][0][]source}{}{}{}}
\makeatother
\title{Three hypotheses under which a relative Proj is quasi-compact}
\author{The Stacks Project Authors}
\date{stacks-project:divisors.tex:local-snapshot-20261012; GNU FDL 1.2 or later}
\begin{document}
\maketitle
\setcounter{subsection}{1}
\noindent\textit{Editorial note.} All delivered mathematical text is the author's own, copied from the pinned Stacks Project chapter file divisors.tex; no statement or proof step was rewritten, reordered or removed.\par\medskip
\noindent\textit{Editorial note.} Carrier window 3: statement and proof both lie inside the window 6000-8000 granted to this lane by the owner lane lane-s13 on 2026-10-14.\par\medskip
\noindent\textit{Editorial note.} The statement carries the underline-Proj notation of the chapter; it is part of the counted statement span and is preserved verbatim.\par\medskip
\noindent\textit{Editorial note.} The references of the counted spans are recorded as non-load-bearing navigation pointers; the argument that consumes them stays inside the counted proof span.\par\medskip
\noindent\textit{Editorial note.} Preamble: sources/stacks-divisors/retained-preamble.tex (88 lines, verbatim present).\par\medskip
\noindent\textit{Allowed standard prerequisites (editorial).}\begin{itemize}
\item Quasi-compactness is local on the base. Schemes, Lemma schemes-lemma-quasi-compact-affine, used at proof lines 7666-7668. Reduces to S affine.
\item Construction of the relative Proj by glueing and its affine model. Constructions, Section constructions-section-relative-proj-via-glueing, used at proof lines 7668-7670. Identifies X with Proj(A).
\item A\_+ is generated by finitely many homogeneous elements. Algebra, Lemma algebra-lemma-S-plus-generated, used at proof lines 7672-7675. Handles case (1).
\item Quasi-compactness of Proj of a graded algebra. Constructions, Lemma constructions-lemma-proj-quasi-compact, used at proof lines 7683-7685. Concludes in both cases.
\end{itemize}\medskip
\ifdefined\hypertarget\hypertarget{sourcepointer-1}{}\fi
\noindent\textit{Source reference.} Original lemma 07Z4: focused-seeds/AGP/sources/stacks-snapshot-20261011/tags.txt, source lines 10134--10134. Navigation pointer into the Stacks Project; the standard fact it supplies is applied inside the counted proof span and the argument resting on it stays there.\par\medskip
\ifdefined\hypertarget\hypertarget{sourcepointer-2}{}\fi
\noindent\textit{Source reference.} Original lemma 01MD: focused-seeds/AGP/sources/stacks-snapshot-20261011/tags.txt, source lines 2118--2118. Navigation pointer into the Stacks Project; the standard fact it supplies is applied inside the counted proof span and the argument resting on it stays there.\par\medskip
\ifdefined\hypertarget\hypertarget{sourcepointer-3}{}\fi
\noindent\textit{Source reference.} Original section 01NM: focused-seeds/AGP/sources/stacks-snapshot-20261011/tags.txt, source lines 2163--2163. Navigation pointer into the Stacks Project; the standard fact it supplies is applied inside the counted proof span and the argument resting on it stays there.\par\medskip
\ifdefined\hypertarget\hypertarget{sourcepointer-4}{}\fi
\noindent\textit{Source reference.} Original lemma 01K4: focused-seeds/AGP/sources/stacks-snapshot-20261011/tags.txt, source lines 2037--2037. Navigation pointer into the Stacks Project; the standard fact it supplies is applied inside the counted proof span and the argument resting on it stays there.\par\medskip
\begin{lemma}
\label{lemma-relative-proj-quasi-compact}
Let $S$ be a scheme. Let $\mathcal{A}$ be a quasi-coherent graded
$\mathcal{O}_S$-algebra. Let
$p : X = \underline{\text{Proj}}_S(\mathcal{A}) \to S$ be the relative
Proj of $\mathcal{A}$. If one of the following holds
\begin{enumerate}
\item $\mathcal{A}$ is of finite type as a sheaf of
$\mathcal{A}_0$-algebras,
\item $\mathcal{A}$ is generated by $\mathcal{A}_1$ as an
$\mathcal{A}_0$-algebra and $\mathcal{A}_1$ is a finite type
$\mathcal{A}_0$-module,
\item there exists a finite type quasi-coherent $\mathcal{A}_0$-submodule
$\mathcal{F} \subset \mathcal{A}_{+}$ such that
$\mathcal{A}_{+}/\mathcal{F}\mathcal{A}$ is a locally nilpotent
sheaf of ideals of $\mathcal{A}/\mathcal{F}\mathcal{A}$,
\end{enumerate}
then $p$ is quasi-compact.
\end{lemma}

\begin{proof}
The question is local on the base, see
Schemes, Lemma \ref{schemes-lemma-quasi-compact-affine}.
Thus we may assume $S$ is affine.
Say $S = \Spec(R)$ and $\mathcal{A}$ corresponds to the
graded $R$-algebra $A$. Then $X = \text{Proj}(A)$, see
Constructions, Section \ref{constructions-section-relative-proj-via-glueing}.
In case (1) we may after possibly localizing more
assume that $A$ is generated by homogeneous elements
$f_1, \ldots, f_n \in A_{+}$ over $A_0$. Then
$A_{+} = (f_1, \ldots, f_n)$ by
Algebra, Lemma \ref{algebra-lemma-S-plus-generated}.
In case (3) we see that $\mathcal{F} = \widetilde{M}$
for some finite type $A_0$-module $M \subset A_{+}$. Say
$M = \sum A_0f_i$. Say $f_i = \sum f_{i, j}$ is the decomposition
into homogeneous pieces. The condition in (3) signifies that
$A_{+} \subset \sqrt{(f_{i, j})}$. Thus in both cases we conclude that
$\text{Proj}(A)$ is quasi-compact by
Constructions, Lemma \ref{constructions-lemma-proj-quasi-compact}.
Finally, (2) follows from (1).
\end{proof}

\end{document}
